% Контрольная 05.10.2026: деревья + плоские графы. Версия 3: три раздела, баллы = трудность задачи.
% Тело для проекта format 179 (Overleaf).
\providecommand{\zb}[1]{\setcounter{lett}{0}\par\refstepcounter{z}\vskip 5pt\noindent\textbf{\arabic{z}}~[#1]~}
\providecommand{\zbb}{\setcounter{lett}{0}\par\refstepcounter{z}\vskip 5pt\noindent\textbf{\arabic{z}}~}
\providecommand{\pb}[1]{\par\vspace{.08cm}\noindent\refstepcounter{lett}\textbf{\asbuk{lett})}~[#1]~}
\providecommand{\razdel}[1]{\par\vspace{9pt}\begin{center}\large\bfseries #1\end{center}\par\vspace{-4pt}}

\ListokName{Контрольная: деревья и плоские графы}
\setcounter{z}{0}

\begin{center}
{\small Школа №\,179. КЛ. 2026/2027 уч. год}\\[3pt]
{\Large\bfseries Контрольная по деревьям и плоским графам}\\[1pt]
{\small 5 октября 2026}
\end{center}

\noindent\textit{Итог подводится по трём задачам, по которым достигнуты наилучшие результаты. Если в задаче несколько пунктов, балл за задачу~--- сумма баллов за пункты. Внутри каждого раздела задачи идут примерно по возрастанию трудности.}

\smallskip
\noindent\textit{Можно пользоваться без доказательства:} для связного плоского графа $\text{В}-\text{Р}+\text{Г}=2$; для плоского графа без петель и кратных рёбер с $\text{В}\ge3$ вершинами $\text{Р}\le3\text{В}-6$.

\razdel{Деревья}

\zb{2} Даня хочет доказать, что утверждения а, б, в, \ldots, я (по одному на каждую букву алфавита, всего~33) равносильны. Какое наименьшее количество импликаций~--- утверждений вида <<из $x$ следует~$y$>>~--- ему придётся доказать?

\zb{3} На плоскости отмечено несколько точек, все попарные расстояния между ними различны. Каждую точку соединили отрезком с ближайшей к ней отмеченной точкой. Может ли среди проведённых отрезков найтись замкнутая ломаная?

\zb{4} Вершины дерева покрашены в красный и синий цвета так, что любые две вершины, соединённые ребром, разного цвета. Вершин больше одной, и красных среди них не меньше, чем синих. Докажите, что какой-то лист дерева красный.

\zbb Скажем, что дерево \textit{нарисовано на двух прямых}, если все его вершины лежат на двух параллельных прямых, каждое ребро~--- отрезок с концами на разных прямых, и рёбра не пересекаются нигде, кроме общих концов. Всего существует 11 попарно неизоморфных деревьев с 7 вершинами.
\pb{2} Нарисуйте на двух прямых как можно больше из них.
\pb{3} Какие из них нарисовать так нельзя? Докажите.

\zb{6} Выпуклый $n$-угольник разбит на треугольники непересекающимися диагоналями, причём в каждой его вершине сходится нечётное число треугольников. Докажите, что $n$ делится на~3.

\razdel{Плоские графы}

\zb{2} Можно ли нарисовать карту из пяти стран, в которой каждые две страны граничат по участку границы (а не только в точке)?

\zbb Шесть домов, каждые два соединены тропинкой. Тропинки могут пересекаться, но в каждой точке пересечения встречаются ровно две тропинки.
\pb{1} Нарисуйте такой план, на котором ровно три точки пересечения.
\pb{2} Докажите, что нельзя убрать две тропинки так, чтобы оставшиеся можно было нарисовать без пересечений.
\pb{2} Докажите, что меньше трёх точек пересечения быть не может.

\zbb Двое по очереди соединяют отрезком две отмеченные точки плоскости. Отрезок не должен пересекать уже проведённые (общие концы можно). Кто не может сделать ход, проиграл. Никакие три отмеченные точки не лежат на одной прямой. Кто выигрывает при правильной игре, если отмечены
\pb{1} три вершины треугольника и одна точка внутри него;
\pb{4} четыре вершины квадрата и пять точек внутри него?

\zbb Петя утверждает: \textit{для любого плоского графа выполнено $\text{В}-\text{Р}+\text{Г}=2$.}

\smallskip
\textit{Доказательство Пети.} Будем стирать рёбра графа по одному. Каждое ребро отделяет одну грань от другой, поэтому, когда его стирают, эти две грани сливаются в одну: Р уменьшается на~1, и Г уменьшается на~1. Значит, $\text{В}-\text{Р}+\text{Г}$ не меняется. Будем стирать рёбра, пока в графе есть циклы. Когда циклов не останется, получится дерево, а у дерева $\text{Р}=\text{В}-1$ и одна грань, так что $\text{В}-\text{Р}+\text{Г}=\text{В}-(\text{В}-1)+1=2$.
\pb{1} Верно ли утверждение Пети?
\pb{2} Найдите все ошибки в доказательстве.
\pb{2} Чему равно $\text{В}-\text{Р}+\text{Г}$ для плоского графа, у которого $k$ компонент связности?

\razdel{Трудные задачи}

\zb{6} В некоторой стране есть 100 городов, которые связаны такой сетью дорог, что из любого города в любой другой можно проехать только одним способом без разворотов. Схема сети дорог известна, развилки и перекрёстки сети необязательно являются городами, всякая тупиковая ветвь сети обязательно заканчивается городом. Навигатор может измерить длину пути по этой сети между любыми двумя городами. Можно ли за 100 таких измерений гарантированно определить суммарную длину всех дорог сети?

\zb{7} Какое наибольшее число клеток доски $9\times9$ можно разрезать по обеим диагоналям, чтобы при этом доска не распалась на несколько частей?

\zb{8} В классе 25 учеников. Каждый заполнил анкету из 25 вопросов, ответив на каждый <<да>> или <<нет>>, и все анкеты различны. Докажите, что можно вычеркнуть один вопрос так, что анкеты останутся попарно различными.

\zb{9} Раскрашенный в чёрный и белый цвета кубик с гранью в одну клетку поставили на одну из клеток шахматной доски и прокатили по ней так, что кубик побывал на каждой клетке ровно по одному разу. Можно ли так раскрасить кубик и так прокатить его по доске, чтобы каждый раз цвета клетки и соприкоснувшейся с ней грани совпадали?

\vfill
\begin{center}\includegraphics[width=.62\textwidth]{Listki/figs/osen.png}\end{center}
